\documentclass[../../../include/open-logic-section]{subfiles}

\begin{document}

\olfileid[fr]{sth}{cardinals}{zfc}
\olsection{\texorpdfstring{$\ZFC$}{ZFC} : un jalon}

Avec l'ajout de l'axiome du bon ordre, nous avons
atteint le dernier jalon théorique. Nous disposons
maintenant de tous les axiomes nécessaires à~$\ZFC$.
Plus précisément :

\begin{defn}
La théorie $\ZFC$ comprend les axiomes suivants :
l'extensionnalité, l'union, la paire, l'ensemble
des parties, l'infini, la fondation, le bon ordre,
ainsi que toutes les instances des schémas de
séparation et de remplacement. Autrement dit,
$\ZFC$ ajoute le bon ordre à~$\ZF$.
\end{defn}

Le sigle $\ZFC$ désigne la théorie des ensembles
de \emph{Zermelo--Fraenkel} avec \emph{choix}.
Cela peut sembler curieux : l'axiome que nous
avons ajouté s'appelait «~bon ordre~», et non
«~choix~». Mais, lorsque nous formulerons le
{choix}, nous verrons que le bon ordre lui est
équivalent (modulo~$\ZF$ ; voir le
\olref[choice][woproblem]{thmwochoice}).
Le choix de notre axiome «~de base~» est donc
indifférent. Quant au nom «~$\ZFC$~», il est
parfaitement standard dans la littérature.

\end{document}
